The $\CauchyEquivalenceSetoidUp$ classifier is the forward lattice edge used when $\RealSeparabilityUp$ reads Cauchy-tail agreement as dense-readback data. It is placed before the stream, regular-sequence, dyadic, and real readback rows are treated as a dense route.

\begin{theorem}[Cauchy equivalence setoid Real-separability forward lattice]
\label{thm:cauchy-equivalence-setoid-real-separability-forward-lattice}
Every $\RealSeparabilityUp$ dense-readback consumer that uses Cauchy-tail agreement must read the displayed $\CauchyEquivalenceSetoidUp$ classifier before the $\StreamNameUp$, $\RegSeqRatUp$, $\DyadicRatCoreUp$, and $\RealUp$ readback rows can be treated as a dense route. The consumer receives a finite classifier edge from the accepted $\CauchyEquivalenceSetoidUp$ packet to \autoref{ch:concrete-instances-real-separability-namecert}; it receives no quotient stream equality, selected representative, countable-choice selector, host real equality, or ambient completeness row.
\end{theorem}
\begin{proof}
Start with the carrier surface of \autoref{def:cauchy-equivalence-setoid-carrier}. Its displayed rows expose the two stream windows, two regular-rational readbacks, dyadic tolerance ledger, setoid-test row, terminal Real seal, componentwise transport, continuation replay, provenance, and local naming data.

The dense-readback consumer may use Cauchy-tail agreement only after the setoid-test row has been inspected. Therefore the $\CauchyEquivalenceSetoidUp$ classifier is the row that turns the surrounding $\StreamNameUp$, $\RegSeqRatUp$, $\DyadicRatCoreUp$, and $\RealUp$ readbacks into a route toward $\RealSeparabilityUp$.

Finally apply the existing Real-separability sibling and route statements, \autoref{thm:cauchy-equivalence-setoid-real-separability-sibling-route} and \autoref{thm:cauchy-equivalence-setoid-real-separability-route}. They place \autoref{ch:concrete-instances-real-separability-namecert} after the same finite classifier boundary, so the forward lattice edge adds no quotient equality, selected representative, countable-choice selector, host real equality, or ambient completeness claim.
\end{proof}
